\chapter{The data pipeline}
\label{ch:pipeline}

A model is only as good as the data it sees. This codebase has an unusually rich input
pipeline: it samples from several datasets, augments aggressively, and even merges a
\emph{second} dataset (for distance) into each training batch. Everything lives under
\file{data\_pipeline/} and is assembled in \file{input\_reader.py}.

\begin{teacher}
In TensorFlow, data loading is built as a \code{tf.data} \emph{pipeline}: a chain of
operations (read $\rightarrow$ decode $\rightarrow$ shuffle $\rightarrow$ augment
$\rightarrow$ batch) that streams examples to the GPU. Each stage runs in parallel on CPU
worker threads while the GPU trains on the previous batch, so data prep and training
overlap.
\end{teacher}

\section{The stage order}
Here is the whole detection pipeline, top to bottom. The order is not arbitrary --- each
stage depends on the previous one.

\begin{center}
\begin{tikzpicture}[node distance=3.6mm,every node/.style={font=\small}]
  \node[blkn] (load) {\code{tfds.load} --- images kept as ENCODED bytes (\code{SkipDecoding})};
  \node[blk,below=of load] (rep) {\code{repeat} each source $\rightarrow$ \code{sample\_from\_datasets} (weights $[95,2,3]$)};
  \node[blk,below=of rep] (shuf) {\code{shuffle} (cheap: KB of encoded bytes, not MB)};
  \node[blk,below=of shuf] (dec) {decode JPEG $\rightarrow$ pre-resize to $672^2$ (uint8)};
  \node[blkg,below=of dec] (cnp) {\textbf{Copy-Paste} (\code{prob=0.2}) --- paste objects \emph{within} an image};
  \node[blkn,below=of cnp] (pb) {\code{padded\_batch(group\_size)} --- polygons padded with $-1.0$};
  \node[blkg,below=of pb] (mos) {\textbf{Mosaic} --- $G$ in $\rightarrow G/R$ out; one \code{random\_perspective} warp};
  \node[blk,below=of mos] (un) {\code{unbatch} $\rightarrow$ \code{shuffle} (disperse a group's outputs)};
  \node[blka,below=of un] (par) {\textbf{Parser} --- polygons $\rightarrow$ radial target; emits \textbf{uint8}};
  \node[blkn,below=of par] (bat) {\code{batch(global\_batch\_size)} $\rightarrow$ \code{prefetch}};
  \node[blka,below=of bat] (gpu) {[on GPU, in \code{train\_step}] HSV jitter + Albumentations + $/255$};
  \foreach \a/\b in {load/rep,rep/shuf,shuf/dec,dec/cnp,cnp/pb,pb/mos,mos/un,un/par,par/bat,bat/gpu}
    \draw[ar] (\a)--(\b);
\end{tikzpicture}
\end{center}

\begin{hood}
Two design choices keep this fast on a CPU-limited host. (1) \textbf{\code{SkipDecoding}}:
images travel through the shuffle buffer as compressed JPEG bytes (kilobytes), and are
only decoded inside the parallel \code{map} --- a 16k-element shuffle buffer stays cheap
instead of holding gigabytes of decoded pixels. (2) \textbf{Color augmentation is moved
to the GPU}: the parsers emit \code{uint8} images and the HSV/Albumentations/$/255$ work
happens once per batch inside \code{train\_step} (\file{data\_pipeline/batch\_color\_aug.py}),
cutting host$\rightarrow$device traffic $4\times$.
\end{hood}

\section{Sampling from several datasets}
The detection stream is not one dataset but several (\code{cleaner\_polygon2026},
\code{field\_misrecog2026}, \code{station\_misrecog}), blended with fixed
\textbf{sampling weights}. Each source is \code{repeat()}-ed first so the blend is a
stationary, infinite stream --- the weights stay constant forever rather than drifting as
a finite source runs out.

\section{Copy-Paste, then Mosaic}
Two augmentations stitch content together, and \emph{order matters}.

\textbf{Copy-Paste} (\file{copy\_paste.py}) pastes a cut-out object (from a separate RGBA
dataset, the alpha channel being the mask) onto the current image. It runs \emph{first},
because it augments \emph{within} a single image. \textbf{Mosaic} then stitches four
whole images together, so it must come after.

\guidefig{fig_perspective.png}{1.0}{The one geometric transform, \code{random\_perspective}
(rotation, scale, shear, translate), rendered by the actual code in
\file{data\_pipeline/augmentations.py}. The box (blue) and polygon (red) ground truth are
transformed by the same matrix and clipped to the image edge. This same warp is applied to
single images and to mosaics.}

\subsection{Mosaic in depth}
Mosaic builds a $2\times$-size canvas, places four images toward a random center, and then
applies \emph{one} \code{random\_perspective} warp that crops it back to $672\times672$.

\guidefig{fig_mosaic.png}{1.0}{The real \code{Mosaic.\_mosaic} output. Four source images
are resized and placed into the four quadrants of a $1344\times1344$ canvas, then a single
warp produces the final image. Ground-truth boxes and polygons are carried through the
\emph{same} matrix, so labels stay perfectly aligned.}

\begin{why}
The mosaic assembles a canvas and warps \emph{once}, rather than warping each of the four
source images separately and compositing. The two are geometrically identical, but the
``warp-once'' canvas form was measured roughly $3\times$ faster on the production CPU
(the warp op is several times more expensive per pixel than a plain resize, and the
alternative pays four full warps per mosaic). On a CPU-bound pipeline this directly sets
the epoch time.
\end{why}

\subsubsection{The group / diversity mechanism}
Mosaic is implemented as a map over a \code{padded\_batch(group\_size)} group. A group of
$G$ decoded images emits $G/R$ output samples, where $R=\code{decodes\_per\_output}$.

\begin{center}
\begin{tikzpicture}[font=\scriptsize,node distance=1.5mm]
  \foreach \i in {0,...,7} {
    \node[blkn,minimum width=7mm,minimum height=6mm] (p\i) at (\i*0.95,0) {\i};
  }
  \node[left=3mm of p0,font=\small] {permutation:};
  % window for output 0 (R=4): 0..3
  \draw[accent,thick,rounded corners] ($(p0.north west)+(-0.05,0.1)$) rectangle ($(p3.south east)+(0.05,-0.1)$);
  \node[accent,above=2mm of p1] {output 0 (4 images)};
  % window for output 1: 4..7
  \draw[green,thick,rounded corners] ($(p4.north west)+(-0.05,0.25)$) rectangle ($(p7.south east)+(0.05,-0.25)$);
  \node[green!50!black,below=2mm of p5] {output 1 (4 images)};
\end{tikzpicture}
\end{center}

\noindent At the default $R=4$ (``stock YOLO''), the width-4 windows \emph{tile} the
permutation: every output uses four distinct images with zero reuse. Lowering $R$ makes
the windows overlap, so each image is reused in $4/R$ outputs with varied partners --- less
decode work, slightly less diversity. Crucially, the number of \emph{training samples} per
epoch is unchanged; $R$ only changes how many source images are decoded.

\begin{pitfall}
The \code{padded\_batch} step pads polygon rows with $-1.0$, \emph{not} the default
$0.0$. Why? $0.0$ is a perfectly valid top-left vertex coordinate --- zero-padded rows
would be read as real vertices and corrupt the polygon target. The value $-1.0$ is the
reserved ``no vertex here'' sentinel throughout the codebase.
\end{pitfall}

\section{Flip and color}
\guidefig{fig_flip_hsv.png}{1.0}{Horizontal flip (real \code{random\_horizontal\_flip} ---
note the box and polygon x-coordinates flip too) and two draws of HSV jitter (real
\code{hsv\_augment}). The brightness term here is \emph{additive}, a deliberate match to the
original codebase. HSV and Albumentations actually run per-batch on the GPU at train time;
this figure calls the same functions directly to show what they do.}

\section{Merging the distance dataset}
Here is the cleverest part of the pipeline. Distance labels come from a \emph{different}
dataset (\code{servingbot\_polygon}) than detection. Rather than build a separate training
job, the two streams are zipped and concatenated on the \textbf{batch dimension}:

\begin{center}
\begin{tikzpicture}[node distance=5mm and 8mm,font=\small]
  \node[blk] (det) {Detection batch\\\scriptsize \code{global\_batch\_size}};
  \node[blkg,below=of det] (dis) {Distance batch\\\scriptsize \code{ignore\_bg=1}};
  \node[blka,right=20mm of $(det)!0.5!(dis)$] (merge) {Merged batch\\\scriptsize concat on axis 0};
  \draw[ar] (det) -- (merge);
  \draw[ar] (dis) -- (merge);
  \node[lbl,above=0.5mm of merge] {\code{tf.data.Dataset.zip}};
\end{tikzpicture}
\end{center}

\noindent The merged batch flows through the same model. Distance-only samples carry
\code{ignore\_bg=1}, which tells the loss to skip class supervision on them (they have no
detection labels). This stream is \textbf{training only} --- there is no distance
validation path, so the distance head is trained but not scored during validation.

\begin{hood}
Three shuffle stages use \emph{distinct} seeds (\code{seed}, \code{seed+1},
\code{seed+2}): the detection source shuffle, the copy-paste source shuffle, and the
post-unbatch shuffle. Distinct seeds keep the stages decorrelated --- a shared seed would
pair the same images every epoch and partially undo each shuffle's job.
\end{hood}
